\documentclass[10pt,a4paper]{amsart}
\usepackage[margin=19mm]{geometry}
\usepackage{amsmath,amssymb,mathtools}
\usepackage[scr=rsfs,cal=euler]{mathalfa}
\usepackage{xfrac}
\usepackage[unicode,hidelinks,bookmarks=false]{hyperref}
\newcommand{\ie}{i.e.\ }
\newcommand{\cf}{cf.\ }
\newcommand{\resp}{respectively\ }
\newcommand{\Subsection}{\subsection}
\providecommand{\nobreakdash}{\nobreak\hbox{-}}
\setlength{\emergencystretch}{2em}
\multlinegap0pt
\usepackage{mathtools}
\usepackage{amssymb}
\usepackage{xcolor}
\usepackage{algorithm}\usepackage{algpseudocode}
\usepackage{enumerate}
\newtheorem{lemma}{Lemma}
\DeclareMathOperator{\dist}{\textup{dist}}
\newcommand{\dlambda}[1]{\,\mathrm d\lambda^{#1}}
\title{Anisotropic Minkowski Content for Countably $\mathcal{H}^k$-rectifiable Sets}
\author{Filip Fry\v{s}}
\date{Source: arXiv:2601.22681v3 (CC BY 4.0)}
\begin{document}
\maketitle

\noindent\emph{Source and licence.} Anisotropic Minkowski Content for Countably $\mathcal{H}^k$-rectifiable Sets. Version arXiv:2601.22681v3 (CC BY 4.0), distributed by arXiv under the Creative Commons Attribution 4.0 International licence, http://creativecommons.org/licenses/by/4.0/, as recorded in the arXiv licence metadata for this exact version and shown on the abstract page https://arxiv.org/abs/2601.22681v3. Full licence notice: \texttt{licenses/arxiv-2601.22681-CC-BY-4.0.txt}.\par\smallskip

\noindent\textit{Editorial scope.} Counted unit: the article's lemma densit (source0.tex lines 606--627). The article's complete original proof of it is delivered with it (source0.tex lines 628--692, one \texttt{proof} environment of 65 lines). No mathematics is written, repaired or re-derived: the statement and the proof are the article's own lines, in the article's own notation, and no retained line is substituted or re-flowed.  No further source-local context is needed: every cross-reference of every retained line resolves inside the two delivered spans.  Nothing was omitted from any retained span, so the package carries no omission record.  No retained line contains a citation, so the wrapper carries no bibliography and no bibliography entry.  No retained line contains a figure environment, an image-inclusion command or drawing code, so no image is delivered and the package declares no asset dependency.  Editorial formatting only, in addition: an AMS \texttt{amsart} preamble that restates the article's own declarations and macros used by the retained lines, namely the environments \texttt{lemma} on separate counters, together with the standard packages those lines need.  The retained environments print in this document as Lemma 1 in order of appearance, whereas the article prints the same items with the numbering of its own full document.  The complete native source of the article as distributed in the arXiv e-print for this exact version is bundled unchanged as \texttt{sources/193659/source0.tex}.
\begin{lemma}\label{densit}
Let $K\subseteq\mathbb{R}^k$ be compact and let $F\in L^1(K)$ be nonnegative.
For every $\varepsilon>0$ and every $\theta\in(0, 1)$ there exist a compact set
$K_{\varepsilon, \theta}\subseteq K$ and a number
$\rho_{\varepsilon,\theta}>0$ such that
\[
\int\limits_{K\setminus K_{\varepsilon,\theta}}F(y)\dlambda{k}(y)<\varepsilon
\]
and
\[
\lambda^k\big(Q_{\mathbb{R}^k}(y, t)\setminus K\big)\leq \theta\,\lambda^k\big(Q_{\mathbb{R}^k}(y, t)\big)
\]
for every $y\in K_{\varepsilon, \theta}$ and every
$0<t<\rho_{\varepsilon, \theta}$.

Moreover, 
\[
Q_{\mathbb{R}^k}(y, t/2)\subseteq K\oplus B_{\mathbb R^k}(0, 3\sqrt{k}\theta^{1/k}t)
\]
for every $y\in K_{\varepsilon, \theta}$ and every
$0<t<\rho_{\varepsilon, \theta}$.
\end{lemma}

\begin{proof}
Since $K$ is compact, it has finite Lebesgue measure. By the Lebesgue density
theorem, for $\lambda^k$-a.e. $y\in K$ we have
\[
\frac{\lambda^k\big(Q_{\mathbb{R}^k}(y, t)\setminus K\big)}{\lambda^k\big(Q_{\mathbb{R}^k}(y, t)\big)}\xrightarrow{t\to0_+}0.
\]
By Egorov's theorem and the absolute continuity of the integral of $F$, there
exists a compact set $K_{\varepsilon, \theta}\subseteq K$ such that
\[
\int\limits_{K\setminus K_{\varepsilon, \theta}}F(y)\dlambda{k}(y)<\varepsilon
\]
and such that the above convergence is uniform on $K_{\varepsilon, \theta}$.
Thus there exists $\rho_{\varepsilon, \theta}>0$ such that
\[
\lambda^k\big(Q_{\mathbb{R}^k}(y, t)\setminus K\big)\leq \theta\,\lambda^k\big(Q_{\mathbb{R}^k}(y, t)\big)
\]
for every $y\in K_{\varepsilon, \theta}$ and every
$0<t<\rho_{\varepsilon, \theta}$.

It remains to prove the inclusion. Let $y\in K_{\varepsilon, \theta}$,
let $0<t<\rho_{\varepsilon, \theta}$, and let $z\in Q_{\mathbb{R}^k}(y, t/2)$. Since
$y\in K$, we have
\[
\dist(z, K)\leq |z-y|\leq \frac{\sqrt{k}}{2}t.
\]
Thus, the desired estimate is trivial if
$\theta^{1/k}\geq 1/6$. We may therefore assume that
$$\theta^{1/k}<1/6.$$

Suppose, to the contrary, that $z\notin K\oplus B_{\mathbb{R}^k}(0, 3\sqrt{k}\theta^{1/k}t).$ Equivalently,
\[
\dist(z, K)>3\sqrt{k}\theta^{1/k}t.
\]
Then the cube
\[
Q_{\mathbb{R}^k}\big(z, 3\theta^{1/k}t/2\big)
\]
is contained in $Q_{\mathbb{R}^k}(y, t)\setminus K$. Indeed, since $z\in Q_{\mathbb{R}^k}(y,t/2)$ and $\theta^{1/k}<1/6$, we have \[ \frac{3}{2}\theta^{1/k}t<\frac{t}{4}, \] and hence \[ Q_{\mathbb{R}^k}\big(z,3\theta^{1/k}t/2\big)\subseteq Q_{\mathbb{R}^k}(y,t). \] Moreover, if $u\in Q_{\mathbb{R}^k}\big(z,3\theta^{1/k}t/2\big)$, then \[ |u-z|\leq \sqrt{k}\,\frac{3}{2}\theta^{1/k}t <3\sqrt{k}\theta^{1/k}t<\dist(z,K), \] and therefore $u\notin K$.


Hence
\[
\lambda^k\big(Q_{\mathbb{R}^k}(y, t)\setminus K\big)
\geq
(3\theta^{1/k}t)^k=\theta (3t)^k.
\]


On the other hand,
\[
\lambda^k\big(Q_{\mathbb{R}^k}(y, t)\setminus K\big)
\leq
\theta\,\lambda^k\big(Q_{\mathbb{R}^k}(y, t)\big)
=
\theta(2t)^k,
\]
which gives a contradiction. Therefore
\[
\dist(z, K)\leq 3\sqrt{k}\theta^{1/k}t
\]
for every $z\in Q_{\mathbb{R}^k}(y, t/2)$, which proves
\[
Q_{\mathbb{R}^k}(y, t/2)\subseteq K\oplus B_{\mathbb R^k}(0, 3\sqrt{k}\theta^{1/k}t).
\]
\end{proof}

\end{document}
